Dans un couloir du cégep, un faucon téméraire installé sur une trottinette électrique fonce à une vitesse de $\FauconVitesse\text{ m/s}$ en direction d'un robot-aspirateur sans défense qui tente d'échapper à l'impact en roulant dans la même direction à une vitesse de $\CibleVitesse\text{ m/s}$. Tout en fonçant, le faucon émet un huissement strident de fréquence $\OiseauFreq\text{ Hz}$. La vitesse du son dans le couloir est de $\DopplerVson\text{ m/s}$.

Un étudiant immobile se tient au centre du couloir, exactement entre le faucon et le robot-aspirateur.
    \begin{enumerate}[label=\alph*)]
        \item Quelle est la fréquence du son direct émis par le faucon et perçu par l'étudiant ?
        \item Quelle est la fréquence de l'écho provenant du robot-aspirateur et perçue par l'étudiant ?
        \item Quel est le nom du phénomène acoustique résultant de la réception simultanée de ces deux sons par l'étudiant, et quelle est la fréquence à laquelle ces moments de faible/forte intensité sonore se répètent ?
    \end{enumerate}

\vspace{0.4cm}